
\newglossaryentry{host}{
  name={host},
  description={
      In the context of virtualization, a host is the physical machine that runs the virtualization software.
      \newline
    }
}

\newglossaryentry{ALU}{
  name={ALU},
  description={
      The Arithmetic Logic Unit is a digital circuit that performs arithmetic and logical operations.
      It is a fundamental building block of the central processing unit \ac{CPU} of a computer.
      \newline
    }
}

\newglossaryentry{Helm chart}
{
  name={Helm chart},
  description={
      Helm filling a similar role a package manager does for an operating system;
      the Helm chart is a collection of files that describe a related set of Kubernetes resources.
      Functioning as a template for a Kubernetes application, Helm charts are used to define, install,
      and upgrade complex Kubernetes applications.
      \newline
    }
}

\newglossaryentry{Node}{
  name={Node},
  description={
      In the context of Clusters a node is a single machine in a cluster.
      It can be a physical machine or a virtual machine.
      \newline
    }
}

\newglossaryentry{zombies}{
  name={zombies},
  description={
      In the context of operating systems, a zombie is a process that has completed execution but still has an entry in the process table.
      It is a process that has finished its execution but has not been removed from the system.
      \newline
    }
}

\newglossaryentry{bus}{
  name={bus},
  description={
      A bus is a communication system that transfers data between components inside a computer or between computers.
      It is a collection of wires that carry data, address, and control signals.
      \newline
    }
}

\newglossaryentry{monkey patching}{
  name={Monkey Patching},
  description={
      Monkey patching is a technique used to modify or extend the behavior of a program at runtime.
      It is a way to change the behavior of a program without changing its source code.
      \newline
    }
}


\newglossaryentry{cordoned}{
  name={cordoned},
  description={
      In the context of Kubernetes, cordoned is a command that marks a node as unschedulable.
      This means that no new pods will be scheduled on the node.
      \newline
    }
}

\newglossaryentry{argsort}{
  name={argsort},
  description={
      Argsort is a function that returns the indices that would sort an array.
      It is a way to get the sorted indices of an array without sorting the array itself.
      \newline
    }
}

\newglossaryentry{drained}{
  name={drained},
  description={
      In the context of Kubernetes, drained is a command that evicts all the pods from a node.
      This is done to prepare the node for maintenance or to remove it from the cluster.
      \newline
    }
}

\newglossaryentry{rootless}{
  name={rootless},
  description={
      Rootless is a security feature that allows a user to run containers without root privileges.
      It is a way to run containers in a more secure way by restricting the capabilities of the container.
      \newline
    }
}

\newglossaryentry{distribution}{
  name={distribution},
  description={
      A distribution is a version of the Linux operating system that is based on the Linux kernel and includes a collection of software applications.
      There are many different distributions of Linux, each with its own set of features and characteristics.
      \newline
    }
}
\newglossaryentry{FPGA}{
  name={FPGA},
  description={
      A FPGA is an integrated circuit device or programmable platform that can be programmed in the field after being manufactured,
      providing electronic product manufacturers with additional design flexibility. Unlike
      application-specific chips, FPGAs allow engineers to make changes very late in the
      design cycle and even upgrade products with new functionality after manufacture.
      \newline
    }
}
\newglossaryentry{guest}{
  name={guest},
  description={
      In the context of virtualization, a guest is anything that is any operating system that is running in a virtual machine.
      on a \gls{host} system.
      \newline
    }
}

\newglossaryentry{root_fs}{
  name={root},
  description={
      The \textbf{root directory} is the top-level directory of the file system hierarchy.
      It is the first directory in the file system and contains all other directories and files.
      To access the root file system, the user must have \gls{root_priv} privileges.
      \newline
    }
}

\newglossaryentry{GitOps}{
  name={GitOps},
  description={
      Git Operations is a way to do Continuous Deployment.
      Its core idea is having a Git repository that always contains declarative descriptions of the infrastructure currently desired in the production environment.
      \newline
    }
}
\newglossaryentry{root_priv}{
  name={root},
  description={
      The \textbf{root user} is the administrative user in a Linux system also known as the superuser or just root.
      The root user has full control over the system and can perform any action on the system.
      It is named after the \gls{root_fs} file system.
      \newline
    }
}

\newglossaryentry{VT-d}{
  name={VT-d},
  description={
      Intel's virtualization technology for directed I/O.
      It allows the virtual machine to access the physical hardware directly.
      \newline
    }
}

\newglossaryentry{virtio}{
  name={virtio},
  description={
      A virtualization standard for network and disk device drivers.
      It is a paravirtualization framework that allows the virtual machine to communicate with the host.
      \newline
    }
}

\newglossaryentry{EPT}{
  name={EPT},
  description={
      Extended Page Tables is a feature of the Intel VT-x virtualization technology.
      It allows the virtual machine to use the physical memory directly.
      \newline
    }
}

\newglossaryentry{kv store}{
  name={key-value store},
  description={
      A key-value store is a type of NoSQL database that uses a simple key-value method to store data.
      It is a simple way to store and retrieve data.
      \newline
    }
}

\newglossaryentry{Single source of truth}{
  name={Single source of truth},
  description={
      Single source of truth is a concept that states that every piece of data should have a single, unambiguous, authoritative representation within a system.
      \newline
    }
}

\newglossaryentry{daemon}{
  name={daemon},
  description={
      A daemon is a background process that runs on a system.
      It is a long-running process that performs a specific task.
      \newline
    }
}

\newglossaryentry{batch system}{
  name={batch system},
  description={
      A batch system is a type of job scheduler that manages and schedules on off batch jobs.
      It is used to run large-scale, long-running, and resource-intensive jobs.
      \newline
    }
}

\newglossaryentry{mico-service}{
  name={microservice},
  description={
      A microservice is a small, independent, and loosely coupled service that is part of a larger application.
      It is designed to be scalable, resilient, and easy to maintain.
      They follow the single responsibility principle and if one service fails irrecoverably, it does not affect the entire application.
    }
}

\newglossaryentry{patching}{
  name={patching},
  description={
      In the context of networking, pactching is the process of connecting
      two or more network segments together on an ISO/OSI layer 1 or 2.
      \newline
    }
}

\newglossaryentry{bridge}{
  name={bridge},
  description={
      A bridge is a network device that connects two or more network segments together on an ISO/OSI layer 2.
      It is used to forward traffic between the segments.
      \newline
    }
}
\newglossaryentry{ISO/OSI model}{
  name={ISO/OSI model},
  description={
      The International Organization for Standardization/Open Systems Interconnection model is a conceptual model that standardizes the functions of a telecommunication or computing system into seven abstraction layers.
      \newline
    }
}

\newglossaryentry{NAT}{
  name={NAT},
  description={
      Network Address Translation is a method of remapping one IP address space into another by modifying network address information in the IP header of packets while they are in transit across a traffic routing device.
      This is used to use a single IP address for multiple devices.
      \newline
    }
}

\newglossaryentry{POSIX}{
  name={POSIX},
  description={
      Portable Operating System Interface is a family of standards specified by the IEEE for maintaining compatibility between operating systems.
      \newline
    }
}

\newglossaryentry{registry}{
  name={registry},
  description={
      In the context of containerization, a registry is a storage and content delivery system that holds container images.
    }
}

\newglossaryentry{HPC}{
  name={High Performance Computing},
  description={
      High Performance Computing is the practice of going beyond the performance a single system can provide, by coordinating multiple systems to work together in parallel.
      The term HPC is used to describe the use of supercomputers and parallel processing techniques for solving complex computational problems.
      The goal of HPC is to provide the necessary computational power to solve problems that are too large or too complex for a single computer to handle.
      HPC is used in a wide range of applications, including scientific research, engineering, weather forecasting, and financial modeling.
      \newline
    }
}

\newglossaryentry{Supercomputing}{
  name={Supercomputing},
  description={
      Supercomputing is the use of supercomputers to solve complex computational problems.
      Supercomputers are high-performance computing systems that are designed to perform large-scale computations at high speeds.
      They are used in a wide range of applications, including scientific research, weather forecasting, and financial modeling.
      \newline
    }
}

\newglossaryentry{bare metal}{
  name={bare metal},
  description={
      Bare metal refers to a computer system that is running directly on the hardware with as few layers of abstraction as possible,
      trying to get a maximum performance out of the system by avoiding the overhead.
      \newline
    }
}